\section{Limitations}
\label{sec:limits}

\begin{enumerate}[leftmargin=*]
  \item \textbf{Small repeat count.} Each configuration cell uses $n{=}3$ repeats; analysis is descriptive only.
  \item \textbf{No inferential testing.} The design does not support significance claims or population generalization.
  \item \textbf{Invalid Qwen LOW cell.} Same-model LOW priming for Qwen 3.8 Max is not estimated.
  \item \textbf{Gemini protocol constraint.} Multi-message Primed delivery differs from single-message packaging used elsewhere.
  \item \textbf{Source length differences.} HIGH $\gg$ LOW $>$ UNSEEN in bytes; magnitude comparability is limited.
  \item \textbf{Genre differences.} HIGH/LOW are narrative; UNSEEN is expository/technical.
  \item \textbf{UNSEEN provenance.} Exact bibliography of the teaching-derived source is unknown.
  \item \textbf{Metric sensitivity.} Canonical coverage can move with lexical/orthographic changes that are not full quality gains.
  \item \textbf{No causal identification.} The paired contrast does not isolate corpus priming from all other session factors.
\end{enumerate}
We do not invent additional limitations beyond those supported by the project artifacts.


These limitations jointly imply that Phase~2B is best read as a carefully documented observational contrast under controlled prompts, not as a definitive assay of Interslavic translation quality or of LLM linguistic competence.
